\chapter{Partial Differentiation}
\label{chap:partial-differentiation}
\lecture{9}{22 May. 00:38}{Ninth Lecture}

\section{Partial Derivatives}

\begin{definition}[Partial Derivative]
  \label{def:partial-derivative}
  Let \( f(x, y) \) be a function of two variables. The \textbf{partial derivative} of \( f \) with respect to \( x \), denoted \( f_x \) or \( \frac{\partial f}{\partial x} \), is found by differentiating \( f \) with respect to \( x \) while keeping \( y \) constant. Similarly, the partial derivative with respect to \( y \), denoted \( f_y \) or \( \frac{\partial f}{\partial y} \), is found by differentiating \( f \) with respect to \( y \) while keeping \( x \) constant.

  \[
    f_x = \frac{\partial f}{\partial x}, \quad
    f_y = \frac{\partial f}{\partial y}
  \]
\end{definition}

\section{First Principle}
\label{sec:first-principle-partial-derivative}

\begin{definition}[First Principle of Partial Derivatives]
  \label{def:partial-derivative-first-principle}
  Let \( f(x, y) \) be a function of two variables. Then the partial derivative of \( f \) with respect to \( x \) at the point \( (x_0, y_0) \) is defined as:

  \[
    f_x(x_0, y_0) = \lim_{h \to 0} \frac{f(x_0 + h, y_0) - f(x_0, y_0)}{h}
  \]

  Similarly, the partial derivative of \( f \) with respect to \( y \) at \( (x_0, y_0) \) is:

  \[
    f_y(x_0, y_0) = \lim_{h \to 0} \frac{f(x_0, y_0 + h) - f(x_0, y_0)}{h}
  \]
\end{definition}

\begin{eg}
  Let \( f(x, y) = x^2y \). Find \( f_x \) and \( f_y \) at \( (1, 2) \) using the first principle.

  For \( f_x \):

  \[
    f_x(1, 2) = \lim_{h \to 0} \frac{f(1 + h, 2) - f(1, 2)}{h}
    = \lim_{h \to 0} \frac{(1 + h)^2 \cdot 2 - 1^2 \cdot 2}{h}
  \]

  \[
    = \lim_{h \to 0} \frac{(1 + 2h + h^2) \cdot 2 - 2}{h}
    = \lim_{h \to 0} \frac{2 + 4h + 2h^2 - 2}{h}
    = \lim_{h \to 0} \frac{4h + 2h^2}{h} = \lim_{h \to 0} (4 + 2h) = 4
  \]

  For \( f_y \):

  \[
    f_y(1, 2) = \lim_{h \to 0} \frac{f(1, 2 + h) - f(1, 2)}{h}
    = \lim_{h \to 0} \frac{1^2 \cdot (2 + h) - 1^2 \cdot 2}{h}
    = \lim_{h \to 0} \frac{2 + h - 2}{h}
    = \lim_{h \to 0} \frac{h}{h} = 1
  \]
\end{eg}

\section{Second Order Partial Derivative}
\label{sec:second-order-partial-derivative}

\begin{definition}[Second Order Partial Derivatives]
  \label{def:second-order-partial-derivative}
  The second order partial derivatives are defined as:
  \[
    f_{xx} = \frac{\partial^2 f}{\partial x^2}, \quad
    f_{yy} = \frac{\partial^2 f}{\partial y^2}, \quad
    f_{xy} = \frac{\partial^2 f}{\partial x \partial y}, \quad
    f_{yx} = \frac{\partial^2 f}{\partial y \partial x}
  \]
  Note: If the mixed partial derivatives \( f_{xy} \) and \( f_{yx} \) are continuous, then \( f_{xy} = f_{yx} \) (\textbf{Clairaut’s Theorem}).
\end{definition}

\begin{eg}
  Let \( f(x, y) = x^2y + 3xy^2 \). Then,
  \[
    f_x = \frac{\partial}{\partial x}(x^2y + 3xy^2) = 2xy + 3y^2
  \]
  \[
    f_y = \frac{\partial}{\partial y}(x^2y + 3xy^2) = x^2 + 6xy
  \]
  \[
    f_{xy} = \frac{\partial}{\partial y}(2xy + 3y^2) = 2x + 6y, \quad f_{yx} = \frac{\partial}{\partial x}(x^2 + 6xy) = 2x + 6y
  \]
\end{eg}

\section{Linearization}

\begin{definition}[Linearization]
  \label{def:linearization-3d}
  The linearization of a function \( f(x, y) \) near a point \( (x_0, y_0) \) is the linear function
  \[
    L(x, y) = f(x_0, y_0) + f_x(x_0, y_0)(x - x_0) + f_y(x_0, y_0)(y - y_0)
  \]
  This is the equation of the tangent plane to the surface \( z = f(x, y) \) at the point \( (x_0, y_0, f(x_0, y_0)) \).
\end{definition}

\begin{eg}
  Let \( f(x, y) = x^2 + y^2 \) and consider the point \( (1, 2) \). Then
  \[
    f(1, 2) = 1^2 + 2^2 = 5,\quad f_x = 2x,\quad f_y = 2y
  \]
  \[
    f_x(1,2) = 2,\quad f_y(1,2) = 4
  \]
  \[
    L(x, y) = 5 + 2(x - 1) + 4(y - 2)
  \]
\end{eg}

\section{Total Differential}

\begin{definition}[Total Differential]
  \label{def:total-diff}
  The total differential \( df \) of a function \( f(x, y) \) is given by:
  \[
    df = f_x(x, y)\,dx + f_y(x, y)\,dy
  \]
  It gives an approximation of the change in \( f \) when \( x \) and \( y \) change by small amounts.
\end{definition}

\begin{eg}
  Let \( f(x, y) = x^2y + y^3 \). Then
  \[
    f_x = 2xy,\quad f_y = x^2 + 3y^2
  \]
  \[
    df = 2xy\,dx + (x^2 + 3y^2)\,dy
  \]
\end{eg}
